% \section{Related Work}

% \paragraph{Insufficient context in RAG.}
% Retrieval failures are a recognized failure point of RAG systems
% \citep{barnett2024seven}. \citet{joren2025sufficient} introduce sufficient
% context as a lens on RAG errors, label contexts as sufficient or insufficient,
% and show that models often answer correctly even from insufficient context and
% hallucinate rather than abstain. They note that a context can narrow the answer
% without settling it, but treat it as insufficient. Models are also known to be
% distracted by irrelevant context \citep{shi2023distracted,yoran2024robust}, and
% \citet{setiawan2026misleading} show that misleading passages collapse
% prompt-based abstention while missing passages do not. Our case differs from
% both: the page is neither irrelevant nor misleading. It is correct and points
% in the right direction, and we show that this alone is enough to weaken
% abstention.

% \paragraph{Partial support and selective answering.}
% Self-RAG \citep{asai2024selfrag} trains a model to judge whether a passage
% fully or partially supports its output, and ALCE \citep{gao2023alce} grades
% citations by whether they support a claim. CRAG \citep{yan2024crag} scores
% retrieved documents and searches again when they are weak, and SURE-RAG
% \citep{qiu2026sure} verifies whether evidence is sufficient before answering.
% These methods recognize partial support, but not what kind of information is
% missing. A family name is a specific, structured kind of partial evidence:
% it says exactly which step is missing and where to look for it.

% \paragraph{Abstention.}
% Surveys and benchmarks document that models often answer when they should
% not \citep{wen2025abstention,kirichenko2025abstentionbench}, including in
% medicine \citep{machcha2026abstain}, and propose methods to detect knowledge
% gaps \citep{feng2024abstain,kadavath2022know}. These works measure whether a
% model declines. We measure which kind of evidence makes it stop declining.

% \paragraph{Ontologies and taxonomies in retrieval.}
% OG-RAG \citep{sharma2024ograg} grounds retrieval in a domain ontology, and
% TaSR-RAG \citep{sun2026tasr} uses a taxonomy to structure multi-hop retrieval.
% The broader program of combining language models with knowledge graphs is
% surveyed by \citet{pan2024unifying}, and neuro-symbolic AI more generally by
% \citet{garcez2023third}. These methods use symbolic structure for every query.
% We study the specific condition under which the page itself signals that a
% symbolic step is needed, and measure how far that step gets.

% \paragraph{What text contains.}
% \citet{mallen2023trust} and \citet{kandpal2023longtail} relate model knowledge
% to entity popularity and to the number of relevant pretraining documents, and
% LMEnt \citep{gottesman2026lment} traces entity mentions through a pretraining
% corpus. These studies work at the level of entities. We measure, for each fact,
% at what level of specificity the text names it.


\section{Related Work}

% \paragraph{Insufficient context in RAG.}
% Retrieval failures are a well-known weakness of RAG pipelines \citep{barnett2024seven}. \citet{joren2025sufficient} examine RAG errors through the lens of context sufficiency, labeling passages as either sufficient or insufficient. They observe that LLMs often hallucinate instead of abstaining when given insufficient context. Although they note that a passage can narrow down the space of possible answers without fully settling the question, their framework treats all such passages simply as insufficient. Prior work also shows that models are easily distracted by irrelevant text \citep{shi2023distracted,yoran2024robust}, and \citet{setiawan2026misleading} find that misleading passages cause prompt-based abstention to fail, whereas completely missing passages do not. Our setting differs from both irrelevant and misleading context. A coarse, family-only passage is factually correct and points in the right direction, yet we show that this partial relevance alone is enough to break abstention.

% \paragraph{Partial support and selective answering.}
% Several methods train or prompt models to evaluate evidence quality. Self-RAG \citep{asai2024selfrag} trains models to judge whether a passage fully or partially supports a generated claim, while ALCE \citep{gao2023alce} evaluates whether citations support model statements. CRAG \citep{yan2024crag} scores retrieved documents and triggers fresh web searches when support is weak, and SURE-RAG \citep{qiu2026sure} verifies context sufficiency before generating an answer. While these systems recognize that support can be partial, they do not identify what specific piece of information is missing. A coarse category mention is a structured form of partial evidence: it indicates that entity-level specificity is missing and points to where the candidates live in an ontology.

% \paragraph{Abstention in language models.}
% Recent surveys and benchmarks show that language models frequently guess when they should abstain \citep{wen2025abstention,kirichenko2025abstentionbench}, including in high-stakes clinical settings \citep{machcha2026abstain}. Prior methods focus on detecting when a model lacks knowledge \citep{kadavath2022know,feng2024abstain}. While those studies measure whether a model abstains overall, we isolate what kind of retrieved evidence tricks a model into abandoning abstention.

% \paragraph{Ontologies and taxonomies in retrieval.}
% Several RAG architectures incorporate structured knowledge. OG-RAG \citep{sharma2024ograg} grounds retrieval in domain ontologies, and TaSR-RAG \citep{sun2026tasr} uses taxonomies to guide multi-hop retrieval. Broader efforts to combine language models with knowledge graphs and symbolic reasoning are surveyed by \citet{pan2024unifying} and \citet{garcez2023third}. Unlike pipelines that invoke symbolic structures on every query, we identify the exact condition where a retrieved text signals the need for an ontology lookup, and we measure where that symbolic handoff succeeds or fails.

% \paragraph{Fact coverage in text corpora.}
% Prior work links a language model's parametric knowledge to how often entities appear in pretraining corpora \citep{mallen2023trust,kandpal2023longtail,gottesman2026lment}. Whereas those studies count mentions of individual entities, we measure the level of granularity at which reference texts state relational facts.


\paragraph{Insufficient context.}
Retrieval failures are a common weakness in RAG systems \citep{barnett2024seven}. \citet{joren2025sufficient} study the errors through the lens of context sufficiency. They label retrieved passages as either sufficient or insufficient and find that LLMs often hallucinate instead of abstaining when context is insufficient. Although they note that a passage can narrow down the answer without fully settling it, they group all such cases under insufficient context. Prior work also shows that models are easily thrown off by irrelevant text \citep{shi2023distracted,yoran2024robust}, and \citet{setiawan2026misleading} show that misleading passages cause prompt-based abstention to fail, whereas missing passages do not. Our focus, \emph{resolvable insufficiency}, differs from both irrelevant and misleading context. A coarse passage is factually accurate and points in the right direction, yet this partial relevance is enough to trigger \clade{} and break abstention.

\paragraph{Partial evidence.}
To handle imperfect passages, several methods train or prompt models to grade evidence quality. Self-RAG \citep{asai2024selfrag} trains a model to judge whether a passage fully or partially supports a generated answer, and ALCE \citep{gao2023alce} checks whether citations back up model claims. Similarly, CRAG \citep{yan2024crag} scores retrieved documents to trigger new searches when support is weak, and SURE-RAG \citep{qiu2026sure} checks whether context is sufficient before answering. While these methods detect that support is partial, they do not identify what kind of information is missing. A coarse mention is a structured form of partial evidence: it shows that fine-grained entity detail is missing and indicates where to find the candidate answers in an ontology.

\paragraph{Abstention.}
When evidence falls short and cannot be resolved, a reliable model should abstain. Recent surveys and benchmarks show that models often guess when they should hold back \citep{wen2025abstention,kirichenko2025abstentionbench}, including in clinical settings \citep{machcha2026abstain}. Most existing methods focus on detecting when a model lacks knowledge so it can decline to answer \citep{kadavath2022know,feng2024abstain}. While those studies measure whether a model abstains overall, we test how the granularity of retrieved evidence (\emph{named}, \emph{family-only}, \emph{silent}) causes a model to stop abstaining and guess.

\paragraph{Ontologies in RAG.}
Rather than guessing or abstaining on a coarse passage, a system can consult a structured ontology. Recent RAG systems use symbolic structures to guide retrieval: OG-RAG \citep{sharma2024ograg} grounds retrieval in domain ontologies, and TaSR-RAG \citep{sun2026tasr} uses taxonomies for multi-hop retrieval. Broader work on combining LLMs with knowledge graphs and symbolic reasoning is surveyed by \citet{pan2024unifying} and \citet{garcez2023third}. Unlike pipelines that query an ontology on every question, we study the specific case where a coarse passage calls for an ontology lookup. We also measure each step of the pipeline: extracting the family, listing its members, and picking the final answer; then we check where the process breaks down.

\paragraph{Fact coverage in text.}
Our analysis also connects to work on what reference and pretraining texts actually state. \citet{mallen2023trust} and \citet{kandpal2023longtail} link a model's closed-book knowledge to entity popularity and document counts, while LMEnt \citep{gottesman2026lment} tracks entity mentions across pretraining data. Whereas those studies count mentions of individual entities, we measure how often reference texts state relational facts at an exact (\emph{named}) level versus a coarse (\emph{family-only}) level.

\paragraph{Fine-grained biology.}
Biology shows the same gap between coarse and fine answers. On Bio-VQA, taxonomic accuracy falls from 0.72 at phylum to 0.12 at species \citep{rayeed2027biovqa}. And models lose accuracy from genus to species \citep{rayeed2025beetleverse} even on curated specimen images with validated traits \citep{rayeed2026beetles}. Knowledge-based VQA benchmarks also contain answers that their corpus does not support \citep{ma2026pitfalls}. We investigate the text side of this gap: a page that names the genus or family but not the answer.
